% !TEX root = ../../main-es.tex
% chapters/es/03-marco-formal.tex
\chapter{Marco formal}
\criterio{plan}{15}
\label{cap:marco-formal}

\porque{Plan de acción / rigor metodológico (EICT 15\%).}{%
  Define el lenguaje formal que los capítulos 4--6 usan. Fallo típico: mezclar sintaxis
  y semántica, o declarar la semántica ``obvia''. Aquí la semántica es la contribución,
  no un trámite.}

\section{Sintaxis del fragmento \dsl{.os}}
\porque{Fragmento formalizado.}{%
  Gramática de: literal, formula, bif, call, ret, ask, ai, map/filter/reduce/foreach,
  fork-join. Tabla de kinds de expresión. Fuente: \texttt{context/features/dsl.md} y
  spec 18. Lo que queda fuera (hot-run, cross-book con targets computados, reflexión)
  se lista explícitamente.}

\section{Semántica denotacional}
\label{sec:semantica}
\porque{Núcleo de O1.}{%
  Función de interpretación $\interp{-}$ del fragmento a morfismos. Totalidad: cada
  término bien tipado denota. Las leyes (asociatividad, identidad) se enuncian aquí y
  se verifican en el cap. \ref{cap:teorema}.}

\section{La SMC del plano fork-join}
\porque{Tensor genuino (O1).}{%
  Especificación 18 ruling 4: CALL multi-target $=$ fork-join sobre snapshots aislados,
  escrituras descartadas, unión determinista en orden de lista. El producto tensorial
  $\tensor$ vale la ley de intercambio — es la base del teorema de no-interferencia.}

\section{La categoría de Freyd del plano secuencial}
\porque{Efectos cuarentenados.}{%
  ADR-003 $=$ plano de control determinista con efecto (LLM) cuarenteno. CALL single-target
  $=$ morfismo no central. La estructura premonoidal/de Freyd separa composición de efectos
  de composición pura.}

\section{Sistema de tipos de referencias}
\porque{Base de O3.}{%
  Tipos de \cellref{}{}: una ref es válida ssi la celda productora publica un esquema que
  habilita el campo. Hoy esto se resuelve en runtime; el verificador (cap.
  \ref{cap:implementacion}) lo sube a compile. Soundness vive en cap. \ref{cap:teorema}.}

% ====== Nomenclatura del marco formal (doble columna, estilo MIT) ======
\begin{nomenclature*}[2em][Nomenclatura del marco formal][section]
\EntryHeading{Categorías y estructuras}
\entry{$\catname{C}$}{categoría base del plano de efectos}
\entry{$\smcc$}{categoría monoidal simétrica del sub-fragmento fork-join}
\entry{$\freyd$}{categoría de Freyd del plano secuencial}
\entry{$\Obj(\catname{C})$}{objetos de $\catname{C}$}
\entry{$\Hom{\catname{C}}{A,B}$}{morfismos $A \to B$ en $\catname{C}$}
\EntryHeading{Constructores}
\entry{$\interp{-}$}{función de interpretación (sintaxis $\to$ morfismos)}
\entry{$\tensor$}{producto tensorial (fork-join paralelo)}
\entry{$\seqcomp$}{composición secuencial ($f \,;\, g$)}
\entry{$\id_A$}{morfismo identidad en $A$}
\EntryHeading{Dominios}
\entry{$\mathrm{Cell}$}{tipo de celda direccionable}
\entry{$\mathrm{Snap}$}{snapshot aislado (input de fork)}
\entry{$\mathrm{LLM}$}{morfismo estocástico (efecto)}
\end{nomenclature*}

